%
% $Header: /home/cvs/root/haskell-report/report/derived.verb,v 1.13 2003/04/08 08:18:19 simonpj Exp $
%
% The paragraph describing the formats of standard representations might
% be deleted, since the info is already in the Prelude.  
% Note that there is a difference in the way readsPrec and showsPrec are defined.
% showsPrec is exact Haskell text, readsPrec uses an auxiliary function which
% isn't quite Haskell.  

%**<title>The Haskell 98 Report: Derived Instances</title>
%**~header

\chapter{Specification of Derived Instances}
\label{derived-appendix}
\index{derived instance}

A {\em derived instance} is an instance declaration that is generated
automatically in conjunction with a \mbox{\tt data} or \mbox{\tt newtype} declaration.
The body of a derived instance declaration is derived syntactically from
the definition of the associated type.  Derived instances are
possible only for classes known to the compiler: those defined in
either the Prelude or a standard library.  In this chapter, we
describe the derivation of classes defined by the Prelude.

If \mbox{$\it T$} is an algebraic datatype declared by:\index{algebraic datatype}
\[\ba{lcl}
\mbox{$\it \makebox{\tt data\ }cx\makebox{\tt \ =>}\ T\ u_1\ \ldots \ u_k$}&\mbox{\tt =}&\mbox{$\it K_1\ t_{11}\ \ldots \ t_{1k_1}\ \makebox{\tt |}\ \cdots\ \makebox{\tt |}\ K_n\ t_{n1}\ \ldots \ t_{nk_n}$}\\
& & \mbox{$\it \makebox{\tt deriving\ (}C_1\makebox{\tt ,}\ \ldots \makebox{\tt ,}\ C_m\makebox{\tt )}$}
\ea\]
(where \mbox{$\it m\geq0$} and the parentheses may be omitted if \mbox{$\it m=1$}) then
a derived instance declaration is possible for a class \mbox{$\it C$} 
if these conditions hold:
\begin{enumerate}
\item
\mbox{$\it C$} is one of \mbox{\tt Eq}, \mbox{\tt Ord}, \mbox{\tt Enum}, \mbox{\tt Bounded}, \mbox{\tt Show}, or \mbox{\tt Read}.

\item
There is a context \mbox{$\it cx'$} such that \mbox{$\it cx'\ \Rightarrow\ C\ t_{ij}$}
holds for each of the constituent types \mbox{$\it t_{ij}$}.

\item
If \mbox{$\it C$} is \mbox{\tt Bounded}, the type must be either an enumeration (all
constructors must be nullary) or have only one constructor.

\item
If \mbox{$\it C$} is \mbox{\tt Enum}, the type must be an enumeration.

\item
There must be no explicit instance declaration elsewhere in the program that
makes \mbox{$\it T\ u_1\ \ldots \ u_k$} an instance of \mbox{$\it C$}.
% or of any of \mbox{$\it C$}'s superclasses.

\item
\hprime{If the data declaration has no constructors (i.e. when \mbox{$\it n=0$}),
then no classes are derivable (i.e. \mbox{$\it m=0$})}
\end{enumerate}
For the purposes of derived instances, a \mbox{\tt newtype} declaration is
treated as a \mbox{\tt data} declaration with a single constructor.

If the \mbox{\tt deriving} form is present,
an instance declaration is automatically generated for \mbox{$\it T\ u_1\ \ldots \ u_k$}
over each class \mbox{$\it C_i$}.
If the derived instance declaration is impossible for any of the \mbox{$\it C_i$}
then a static error results.
If no derived instances are required, the \mbox{\tt deriving} form may be
omitted or the form \mbox{\tt deriving\ ()} may be used.

% OLD:
%If the \mbox{\tt deriving} form is omitted then instance
%declarations are derived for the datatype in as many of the six
%classes mentioned above as is possible; that is, no static error can occur.
%Since datatypes may be recursive, the implied inclusion in
%these classes may also be recursive, and the largest
%possible set of derived instances is generated.  For example,
%\bprog
%@%@
%data  T1 a = C1 (T2 a) | Nil1
%data  T2 a = C2 (T1 a) | Nil2
%@%@
%\eprog
%Because the \mbox{\tt deriving} form is omitted, we would expect derived
%instances for \mbox{\tt Eq} (for example).  But \mbox{\tt T1} is in \mbox{\tt Eq} only if \mbox{\tt T2}
%is, and \mbox{\tt T2} is in \mbox{\tt Eq} only if \mbox{\tt T1} is.  The largest solution has
%both types in \mbox{\tt Eq}, and thus both derived instances are generated.

Each derived instance declaration will have the form:
\[
\mbox{$\it \makebox{\tt instance\ (}cx\makebox{\tt ,\ }cx'\makebox{\tt )\ =>}\ C_i\ (T\ u_1\ \ldots \ u_k)\ \makebox{\tt where}\ \makebox{\tt {\char'173}}\ d\ \makebox{\tt {\char'175}}$}
\]
where \mbox{$\it d$} is derived automatically depending on \mbox{$\it C_i$} and the data
type declaration for \mbox{$\it T$} (as will be described in the remainder of
this section).
%% Yale nuked this:
%% The class assertions \mbox{$\it C'\ u'$} are constraints on \mbox{$\it T$}'s
%% type variables that are inferred from the instance declarations of the
%% constituent types \mbox{$\it t_{ij}$}.  For example, consider:
%% \bprog
%% @
%% data  T1 a = C1 (T2 a) deriving Eq
%% data  T2 a = C2 a      deriving ()
%% @
%% \eprog
%% And consider these three different instances for \mbox{\tt T2} in \mbox{\tt Eq}:\nopagebreak
%% \bprog
%% @
%% instance            Eq (T2 a) where (C2 x) == (C2 y)  =  True
%% 
%% instance (Eq  a) => Eq (T2 a) where (C2 x) == (C2 y)  =  x == y
%% 
%% instance (Ord a) => Eq (T2 a) where (C2 x) == (C2 y)  =  x > y
%% @
%% \eprog
%% The corresponding derived instances for \mbox{\tt T1} in \mbox{\tt Eq} are:
%% \bprog
%% @
%% instance            Eq (T1 a) where (C1 x) == (C1 y)  =  x == y
%% 
%% instance (Eq  a) => Eq (T1 a) where (C1 x) == (C1 y)  =  x == y
%% 
%% instance (Ord a) => Eq (T1 a) where (C1 x) == (C1 y)  =  x == y
%% @
%% \eprog

The context \mbox{$\it cx'$} is the smallest context satisfying point (2) above.
For mutually recusive data types, the compiler may need to perform a
fixpoint calculation to compute it.

The remaining details of the derived
instances for each of the derivable Prelude classes are now given.
Free variables and constructors used in these translations
always refer to entities defined by the \mbox{\tt Prelude}.

\section{Derived instances of \texorpdfstring{\mbox{\tt Eq}}{Eq} and \texorpdfstring{\mbox{\tt Ord}}{Ord}}
\indexdi{Eq}
\indexdi{Ord}
The class methods automatically introduced by derived instances
of \mbox{\tt Eq} and \mbox{\tt Ord} are \mbox{\tt (==)},
\indextt{==}
\mbox{\tt (/=)},
\indextt{/=}
\mbox{\tt compare}\indextt{compare},
\mbox{\tt (<)},
\indextt{<}
\mbox{\tt (<=)},
\indextt{<=}
\mbox{\tt (>)},
\indextt{>}
\mbox{\tt (>=)},
\indextt{>=}
\mbox{\tt max}\indextt{max}, and 
\mbox{\tt min}\indextt{min}.  The latter seven operators are defined so
as to compare their arguments lexicographically with respect to
the constructor set given, with earlier constructors in the datatype
declaration counting as smaller than later ones.  For example, for the
\mbox{\tt Bool} datatype, we have that \mbox{\tt (True\ >\ False)\ ==\ True}.

Derived comparisons always traverse constructors from left to right.
These examples illustrate this property:
\begin{quote}
\mbox{\tt (1,undefined)\ ==\ (2,undefined)\ }\mbox{$\it \Rightarrow$}\mbox{\tt \ \ \ \ False}\\
\mbox{\tt (undefined,1)\ ==\ (undefined,2)\ }\mbox{$\it \Rightarrow$}\mbox{\tt \ \ \ \ }\mbox{$\it \bot$}
\end{quote}
All derived operations of class \mbox{\tt Eq} and \mbox{\tt Ord} are strict in both arguments.
For example, \mbox{$\it \makebox{\tt False\ <=\ }\bot$} is \mbox{$\it \bot$}, even though \mbox{\tt False} is the first constructor
of the \mbox{\tt Bool} type.

\section{Derived instances of \texorpdfstring{\mbox{\tt Enum}}{Enum}}
\indexdi{Enum}
Derived instance declarations for the class \mbox{\tt Enum} are only
possible for enumerations (data types with only nullary constructors).

The nullary constructors are assumed to be
numbered left-to-right with the indices 0 through $n-1\/$.
The \mbox{\tt succ} and \mbox{\tt pred} operators give the successor and predecessor
respectively of a value, under this numbering scheme.  It is
an error to apply \mbox{\tt succ} to the maximum element, or \mbox{\tt pred} to the minimum
element.

The \mbox{\tt toEnum} and \mbox{\tt fromEnum} operators map enumerated values to and
from the \mbox{\tt Int} type; \mbox{\tt toEnum} raises a runtime error if the \mbox{\tt Int} argument
is not the index of one of the constructors.

The definitions of the remaining methods are 
\par
{\small
\bprog
{\small\begin{verbatim}
  enumFrom x           = enumFromTo x lastCon
  enumFromThen x y     = enumFromThenTo x y bound
                       where
                         bound | fromEnum y >= fromEnum x = lastCon
                               | otherwise                = firstCon
  enumFromTo x y       = map toEnum [fromEnum x .. fromEnum y]
  enumFromThenTo x y z = map toEnum [fromEnum x, fromEnum y .. fromEnum z]
\end{verbatim}}
\eprog
}
where \mbox{\tt firstCon} and \mbox{\tt lastCon} are respectively the first and last
constructors listed in the \mbox{\tt data} declaration.
For example,
given the datatype:
\bprog
{\small\begin{verbatim}
  data  Color = Red | Orange | Yellow | Green  deriving (Enum)
\end{verbatim}}
\eprog
we would have:
\bprog
{\small\begin{verbatim}
  [Orange ..]         ==  [Orange, Yellow, Green]
  fromEnum Yellow     ==  2
\end{verbatim}}
\eprog

\section{Derived instances of \texorpdfstring{\mbox{\tt Bounded}}{Bounded}}
\indexdi{Bounded}
The \mbox{\tt Bounded} class introduces the class
methods 
\mbox{\tt minBound}\indextt{maxBound} and
\mbox{\tt maxBound}\indextt{minBound},
which define the minimal and maximal elements of the type.
For an enumeration, the first and last constructors listed in the
\mbox{\tt data} declaration are the bounds.  For a type with a single
constructor, the constructor is applied to the bounds for the
constituent types.  For example, the following datatype:
\bprog
{\small\begin{verbatim}
  data  Pair a b = Pair a b deriving Bounded
\end{verbatim}}
\eprog
would generate the following \mbox{\tt Bounded} instance:
\bprog
{\small\begin{verbatim}
  instance (Bounded a,Bounded b) => Bounded (Pair a b) where
    minBound = Pair minBound minBound
    maxBound = Pair maxBound maxBound
\end{verbatim}}
\eprog

\section{Derived instances of \texorpdfstring{\mbox{\tt Read}}{Read} and \texorpdfstring{\mbox{\tt Show}}{Show}}
\label{derived-text}
\indexdi{Read}
\indexdi{Show}
The class methods automatically introduced by derived instances
of \mbox{\tt Read} and \mbox{\tt Show} are \mbox{\tt showsPrec}\indextt{showsPrec},
\mbox{\tt readsPrec}\indextt{readsPrec},
\mbox{\tt showList}\indextt{showList}, and \mbox{\tt readList}\indextt{readList}.
They are used to coerce values into strings and parse strings into
values.

The function \mbox{\tt showsPrec\ d\ x\ r} accepts a precedence level \mbox{\tt d}
(a number from \mbox{\tt 0} to \mbox{\tt 11}), a value \mbox{\tt x}, and a string \mbox{\tt r}.
It returns a string representing \mbox{\tt x} concatenated to \mbox{\tt r}.
\mbox{\tt showsPrec} satisfies the law:
\[
\mbox{$\it \makebox{\tt showsPrec\ d\ x\ r\ ++\ s\ \ ==\ \ showsPrec\ d\ x\ (r\ ++\ s)}$}
\]
The representation will be enclosed in parentheses if the precedence
of the top-level constructor in \mbox{\tt x} is less than \mbox{\tt d}.  Thus,
if \mbox{\tt d} is \mbox{\tt 0} then the result is never surrounded in parentheses; if
\mbox{\tt d} is \mbox{\tt 11} it is always surrounded in parentheses, unless it is an
atomic expression (recall that function application has precedence \mbox{\tt 10}).
The extra parameter \mbox{\tt r} is essential if tree-like
structures are to be printed in linear time rather than time quadratic
in the size of the tree.

The function \mbox{\tt readsPrec\ d\ s} accepts a precedence level \mbox{\tt d} (a number
from \mbox{\tt 0} to \mbox{\tt 10}) and a string \mbox{\tt s}, and attempts to parse a value 
from the front of the string, returning a list of (parsed value, remaining string) pairs.
If there is no successful parse, the returned list is empty.
Parsing of an un-parenthesised infix operator application succeeds only 
if the precedence of the operator is greater than or equal to \mbox{\tt d}.

It should be the case that
\begin{center}
\mbox{\tt (x,"")}  is an element of  \mbox{\tt (readsPrec\ d\ (showsPrec\ d\ x\ ""))}
\end{center}
That is, \mbox{\tt readsPrec} should be able to parse the string produced
by \mbox{\tt showsPrec}, and should deliver the value that \mbox{\tt showsPrec} started
with.

\mbox{\tt showList} and \mbox{\tt readList} allow lists of objects to be represented
using non-standard denotations.  This is especially useful for strings
(lists of \mbox{\tt Char}).

%Because a string is a list of characters, \mbox{\tt showsPrec} of a string
%such as \mbox{\tt "abc"} will result in the string 
%\mbox{\tt "[}\fwq\mbox{\tt a}\fwq \mbox{\tt ,} \fwq\mbox{\tt b}\fwq \mbox{\tt ,} \fwq\mbox{\tt c}\fwq \mbox{\tt ]"}.  Because
%\mbox{\tt "{\char'134}"abc{\char'134}""} would be a better representation,
%the operators \mbox{\tt showList}
%and \mbox{\tt readList} are provided in the class \mbox{\tt Text} for coercing {\em
%lists} of values to and from strings.  In particular, \mbox{\tt showsPrec} of a
%string will yield the verbose form above, and \mbox{\tt showList} will yield
%the compact version.  For most other datatypes, \mbox{\tt showList} and
%\mbox{\tt readList} will yield the same result as \mbox{\tt showsPrec} and \mbox{\tt readsPrec}.


%The instances of \mbox{\tt Text} for the standard types \mbox{\tt Int}, \mbox{\tt Integer}, \mbox{\tt Float},
%\mbox{\tt Double}, \mbox{\tt Char}, lists, tuples, and rational and complex numbers are
%defined in the  
%standard prelude (see Chapter~\ref{stdprelude}).
%For characters and strings, the control characters that have special
%representations (\mbox{\tt {\char'134}n} etc.) are shown as such by \mbox{\tt showsPrec};
%otherwise, ASCII mnemonics are used.
%Non-ASCII characters are shown by decimal escapes.
%Floating point numbers are represented by decimal numbers
%of sufficient precision to guarantee \mbox{\tt read\ .\ show} is an identity
%function.  If $b$ is the floating-point radix and there are
%$w$ base-$b$ digits in the floating-point significand,
%the number of decimal digits required is
%$d = \lceil w \log_{10} b \rceil + 1$ \cite{matula70}.
%Numbers are shown in non-exponential format if this requires
%only $d$ digits; otherwise, they are shown in exponential format,
%with one digit before the decimal point.  \mbox{\tt readsPrec} allows
%an exponent to be unsigned or signed with \mbox{\tt +} or \mbox{\tt -}; \mbox{\tt showsPrec}
%shows a positive exponent without a sign.

\mbox{\tt readsPrec} will parse any valid representation of the standard types 
apart from strings, for which only quoted strings are accepted, and other lists,
for which only the bracketed form \mbox{\tt [}\ldots\mbox{\tt ]} is accepted. See
Chapter~\ref{stdprelude} for full details.

The result of \mbox{\tt show} is a syntactically correct \Haskell{} expression
containing only constants, given the fixity declarations in force at
the point where the type is declared.  It contains only the
constructor names defined in the data type, parentheses, and
spaces. When labelled constructor fields are used, braces, commas,
field names, and equal signs are also used.  Parentheses
are only added where needed, {\em ignoring associativity}.  No line breaks
are added. The result of \mbox{\tt show} is readable by \mbox{\tt read} if all component
types are readable.  (This is true for all instances defined in the
Prelude but may not be true for user-defined instances.)

Derived instances of \mbox{\tt Read} make the following assumptions, 
which derived instances of \mbox{\tt Show} obey:
\begin{itemize}
\item If the constructor is defined to be an infix operator, then 
the derived \mbox{\tt Read} instance will parse only infix applications of the 
constructor (not the prefix form).

\item Associativity is not used to reduce the occurrence of 
parentheses, although precedence may be. For example, given
\bprog
{\small\begin{verbatim}
  infixr 4 :$
  data T = Int :$ T  |  NT
\end{verbatim}}
\eprog
then:
\begin{itemize}
\item \mbox{\tt show\ (1\ :{\char'44}\ 2\ :{\char'44}\ NT)} produces the string \mbox{\tt "1\ :{\char'44}\ (2\ :{\char'44}\ NT)"}.
\item \mbox{\tt read\ "1\ :{\char'44}\ (2\ :{\char'44}\ NT)"} succeeds, with the obvious result.
\item \mbox{\tt read\ "1\ :{\char'44}\ 2\ :{\char'44}\ NT"} fails.
\end{itemize}

\item 
If the constructor is defined using record syntax, the derived \mbox{\tt Read} 
will parse only the record-syntax form, and furthermore, the fields must be 
given in the same order as the original declaration.

\item The derived \mbox{\tt Read} instance allows arbitrary Haskell whitespace between 
tokens of the input string.  Extra parentheses are also allowed.
\end{itemize}


% However, the
% derived \mbox{\tt Read} and \mbox{\tt Show} instances have the following properties:
% \begin{itemize}
% \item The result of \mbox{\tt show} is a syntactically correct \Haskell{}
%       expression containing only constants
%       given the fixity declarations in force at the point where
%       the type is declared.
% \item The result of \mbox{\tt show} is readable by \mbox{\tt read} if all component
%       types are readable.  (This is true for all instances defined in
%       the Prelude but may not be true for user-defined instances.)
% \item The instance generated by \mbox{\tt Read} allows arbitrary whitespace
%       between tokens on the input string.  Extra parentheses are also 
%       allowed.
% \item The result of \mbox{\tt show} contains only the constructor names defined
%       in the data type, parentheses, and spaces.  When labelled
%       constructor fields are used, braces, commas, field names, and
%       equal signs are also used.
%       Spaces and parentheses are only added where
%       needed.  No line breaks are added.
% \item If a constructor is defined using labelled field syntax then the derived
%       \mbox{\tt show} for that constructor will use this same
%       syntax; the fields will be in the order declared in the
%       \mbox{\tt data} declaration.  The derived \mbox{\tt Read} instance will use
%       this same syntax: all fields must be present and the declared order
%       must be maintained.
% \item If a constructor is defined in the infix style, the derived \mbox{\tt Show}
%       instance will also use infix style.  The derived \mbox{\tt Read} instance will
%       require that the constructor be infix.
% \end{itemize}

The derived \mbox{\tt Read} and \mbox{\tt Show} instances may be unsuitable for some
uses.  Some problems include:
\begin{itemize}
\item Circular structures cannot be printed or read by these
instances.
\item The printer loses shared substructure; the printed
representation of an object may be much larger than necessary.
\item The parsing techniques used by the reader are very inefficient;
reading a large structure may be quite slow.
\item There is no user control over the printing of types defined in
the Prelude.  For example, there is no way to change the
formatting of floating point numbers.
\end{itemize}


%Figure~\ref{derived-text} gives the general form of a derived instance
%of \mbox{\tt Text} for a user-defined datatype:
%\[
%\mbox{$\it \makebox{\tt data}\ cx\ \makebox{\tt =>}\ T\ u_1\ \ldots \ u_k\ \makebox{\tt =}\ \ldots $}
%\]
%\mbox{\tt showsPrec} and \mbox{\tt readsPrec} are as
%in Appendices~\ref{showsPrec-spec} and \ref{readsPrec-spec}.  The omitted
%definitions of \mbox{\tt readList} and \mbox{\tt showList} in
%Figure~\ref{standard-classes} (page~\pageref{standard-classes})
%are:
%\bprog
%@
%readList:: ReadS [a]
%readList r = [pr | (\mbox{$\it [$},s) <- lex r,
%                 pr      <- readl s    ]
%           where readl s = [([],t) | (\mbox{$\it ]$},t) <- lex s] ++
%                          [(x:xs,v) |  (x,t) <- reads s,
%                                       (\mbox{$\it ,$},u) <- lex t,
%                                       (xs,v) <- readl u       ]
%
%showList:: [a] -> ShowS
%showList xs = showChar '[' . showl xs
%            where
%            showl [] = showChar ']'
%            showl (x:xs) = shows x . showChar ',' . showl xs
%@
%\eprog
%\begin{figure}
%\outline{
%@
%instance (C, Text u1, ..., Text uk) => Text (T u1 ... uk) where
%       showsPrec = ...
%       readsPrec = ...
%@
%}
%\caption{General form of a derived instance of \mbox{\tt Text}}
%\label{derived-text}
%\end{figure}


\section{An Example}

As a complete example, consider a tree datatype:\nopagebreak
%\bprog
%@
%data Tree a = Leaf a | Tree a :^: Tree a
%@
%\eprog\nopagebreak
%Since there is no \mbox{\tt deriving} clause, this is shorthand for:\nopagebreak
\bprog
{\small\begin{verbatim}
  data Tree a = Leaf a | Tree a :^: Tree a
       deriving (Eq, Ord, Read, Show)
\end{verbatim}}
\eprog
Automatic derivation of
instance
declarations for \mbox{\tt Bounded} and \mbox{\tt Enum} are not possible, as \mbox{\tt Tree} is not
an enumeration or single-constructor datatype.  The complete
instance declarations for \mbox{\tt Tree} are shown in Figure~\ref{tree-inst},
Note the implicit use of default class method
\index{default class method}
definitions---for
example, only \mbox{\tt <=} is defined for \mbox{\tt Ord}, with the other
class methods (\mbox{\tt <}, \mbox{\tt >}, \mbox{\tt >=}, \mbox{\tt max}, and \mbox{\tt min}) being defined by the defaults given in
the class declaration shown in Figure~\ref{standard-classes}.

\begin{figure}[tb]
\begin{outlineenv}
\small
{\small\begin{verbatim}
infixr 5 :^:
data Tree a =  Leaf a  |  Tree a :^: Tree a

instance (Eq a) => Eq (Tree a) where
        Leaf m == Leaf n  =  m==n
        u:^:v  == x:^:y   =  u==x && v==y
             _ == _       =  False

instance (Ord a) => Ord (Tree a) where
        Leaf m <= Leaf n  =  m<=n
        Leaf m <= x:^:y   =  True
        u:^:v  <= Leaf n  =  False
        u:^:v  <= x:^:y   =  u<x || u==x && v<=y

instance (Show a) => Show (Tree a) where

        showsPrec d (Leaf m) = showParen (d > app_prec) showStr
          where
             showStr = showString "Leaf " . showsPrec (app_prec+1) m

        showsPrec d (u :^: v) = showParen (d > up_prec) showStr
          where
             showStr = showsPrec (up_prec+1) u . 
                       showString " :^: "      .
                       showsPrec (up_prec+1) v
                -- Note: right-associativity of :^: ignored

instance (Read a) => Read (Tree a) where

        readsPrec d r =  readParen (d > up_prec)
                         (\r -> [(u:^:v,w) |
                                 (u,s) <- readsPrec (up_prec+1) r,
                                 (":^:",t) <- lex s,
                                 (v,w) <- readsPrec (up_prec+1) t]) r

                      ++ readParen (d > app_prec)
                         (\r -> [(Leaf m,t) |
                                 ("Leaf",s) <- lex r,
                                 (m,t) <- readsPrec (app_prec+1) s]) r

up_prec  = 5    -- Precedence of :^:
app_prec = 10   -- Application has precedence one more than
                -- the most tightly-binding operator
\end{verbatim}}
\end{outlineenv}
%**<div align=center> <h4>Figure 8</h4> </div>
\ecaption{Example of Derived Instances}
\label{tree-inst}
\end{figure}

%**~footer
